\chapter{Proof of Stake}\label{chapter.stake}

In \Cref{chapter.untrusty-world}, we established that money is a social construct:
Money is made valuable by a counter-party willing to accept it and hold it;
\emph{society at large} must be willing to accept it and hold it.
As a holder of money, we trust that the majority of society will act in this manner,
or, at the very least, those members of society who are economically active,
the \emph{economic majority}. It is only natural to assume that the majority of
coins in our system are in honest hands.

This is what gave rise to the assumption we made when we designed proof-of-work consensus,
namely that the majority of computational power is in honest hands. Since the majority of
money belongs to honest parties, and these parties are free to purchase computational power,
we deduce that the majority of computational power is honest.
% This conclusion is imperfect...

{\color{red}
\begin{itemize}
\item The proof-of-work environmental impact
\item The honest stake assumption
\item Grinding attacks
\item Proof-of-stake time: The slot
\item Combining honest and adversarial randomness to obtain unbiased randomness
\item Pseudorandomness: Getting more randomness from little randomness
\item Epoch randomness
\item Slot leaders
\item The Ouroboros protocol
\item k+1 blocks in any continuous 2k slots
\item Secret Sharing Schemes
\item Verifiable Secret Sharing
\item Publicly Verifiable Secret Sharing
\item Force-opening commitments using honest majority
\item Verifiable random functions
\item The proof-of-stake inequality
\item The Ouroboros Praos protocol
\item Aggregation independence
\end{itemize}
}
